\section{Future Work}
The contact equivalences reduce prime and twin-prime selection to exact
levels of continuous functions. Quantitative information still requires
control of the integer samples. For a finite set $I$ of integer starts
$n\ge2$, the binary values from Section~11 give the exact candidate count
\begin{equation*}
\begin{aligned}
N_M(I)&=\sum_{n\in I}A_M(n).
\end{aligned}
\end{equation*}

Every summand is one for a surviving pair and zero otherwise. A positive
independent lower bound in a region where both endpoints are certified would
therefore produce a twin pair. Direct evaluation of the sum is an exact
count, while proving such a bound without first enumerating its contacts is
the remaining analytic question.

A continuous integral requires a separate sampling argument before it can
control this count. Positive area between integers does not by itself force
a unit-amplitude integer sample. The bounded product is useful for separating
height magnitude from contact selection, but its boundedness alone gives no
lower bound for $N_M(I)$.

The folding identity suggests investigating which translated wave groups can
be combined while keeping their activation defects. Quadratic arguments such
as $n(n+2)$ and coordinates restricted to $6k-1,6k+1$ suggest further
arithmetic organizations of the pair condition. Estimates arising from these
choices remain research questions here; no new distribution result follows
from the change of variables.

There is also a separate coordinate comparison. A cycle integral enumerates
surviving values by position in an ordered sequence. The curves here use the
numerical value as their real coordinate. Turning a contact set into that
ordered sequence still requires enumeration and reindexing. Finite-stage
tail periodicity and the Fourier description of a fixed gap cycle are not
developed in this article. Their comparison with the activated curves is a
structural extension, not an assertion that continuity has replaced the
survivor-index construction.
